\section{Formal development}\label{app:lean}

The correctness of the solvers rests on short combinatorial arguments and on
exact formulas that the implementation evaluates millions of times per
second.  An error in a sign or a boundary case would not crash the code; it
would silently produce worse or invalid results.  We formalised part of the
mathematics in Lean~4~\cite{deMouraU21lean4} with
Mathlib~\cite{mathlib2020}.  The development (\texttt{lean/CCProofs}, about
1{,}800 lines in fourteen files) is checked by the Lean kernel without
\texttt{sorry} and depends only on the axioms \texttt{propext},
\texttt{Classical.choice} and \texttt{Quot.sound}.  Table~\ref{tab:lean} lists
what it covers.  It does \emph{not} cover Theorem~\ref{thm:sparse}, the block
step and the convergence statement of Theorem~\ref{thm:anytime}(b,c), the gap
decomposition (Theorem~\ref{thm:gap}) or its local corollary, the $2.06$ part
of Theorem~\ref{thm:certiflip}, or any code.  The Pivot guarantee enters as a
hypothesis.

The formal model is that of Section~\ref{sec:prelim}: a finite vertex set, a
simple graph $G^+$ and a labelling $c:V\to\alpha$ into a type with decidable
equality.  The cost is summed over ordered pairs, so every count appears
multiplied by two (\texttt{cost\_eq\_two\_mul\_card} relates it to the
unordered count).  The weighted instance has integer node sizes and symmetric
weights, as in~\eqref{eq:wcost}.

\begin{table}[!tbp]
\centering
\caption{Coverage of the formal development (Lean~4, Mathlib).  ``Full'': the statement of the
paper is proved as stated (up to the ordered-pair convention); ``partial'':
the part named is proved; ``reduction'': proved from the cited guarantee
taken as a hypothesis; ``--'': not formalised.}\label{tab:lean}
\small
\begin{tabular}{p{0.30\linewidth}p{0.12\linewidth}p{0.46\linewidth}}
\toprule
result & coverage & Lean theorems \\
\midrule
Lemma~\ref{lem:half} & partial & far pairs are separated by every optimal clustering: \texttt{optimal\_separates\_far}, \texttt{sepFar\_of\_optimal} \\
Theorem~\ref{thm:sparse} & -- & \\
Proposition~\ref{prop:subgraph}(i) & full & \texttt{star\_row\_valid} (star rows in $x$-form, far-separating clusterings) \\
Proposition~\ref{prop:subgraph}(ii) & partial & bad triangles and induced stars: \texttt{bad\_triangle\_costIn}, \texttt{star\_bound}; packings: \texttt{packing\_bound} \\
rows \eqref{eq:t1}--\eqref{eq:t2} & full & \texttt{sep\_triangle}, \texttt{sep\_far\_triangle} \\
Theorem~\ref{thm:anytime}(a) & full & \texttt{far\_dual\_bound}, \texttt{far\_dual\_le\_opt} (bound on $P$ for rows valid for far-separating clusterings, $\le$ the cost of every clustering); \texttt{cost\_eq\_supp} \\
Theorem~\ref{thm:anytime}(b,c), Theorem~\ref{thm:gap}, Corollary~\ref{cor:local} & -- & \\
Theorem~\ref{thm:certiflip} & (a) reduction, (b) full, (c) -- & (a) as Theorem~\ref{thm:pxmem}; (b) is Theorem~\ref{thm:anytime}(a) \\
Proposition~\ref{prop:checker} & partial & the inequality evaluated by the checker: \texttt{far\_dual\_bound}, \texttt{far\_dual\_le\_opt}; the checker code itself is not verified \\
Lemma~\ref{lem:twins} & full & \texttt{twins\_split\_improvable}, \texttt{optimal\_keeps\_twins} \\
Pivot keeps twins together & partial & one round: \texttt{pivot\_round\_twins} \\
contraction~\eqref{eq:wcost} & partial & cost identity: \texttt{cost\_contract} \\
move and swap formulas~\eqref{eq:move} & full & \texttt{cost\_move}, \texttt{costW\_move}, \texttt{costW\_swap}, \texttt{cost\_update} \\
Theorem~\ref{thm:px} & full & any pairwise-additive cost: \texttt{crossoverP\_le}; unweighted: \texttt{crossover\_le}; weighted: \texttt{crossoverW\_le}; relabelling: \texttt{cost\_comp\_injective}, \texttt{costW\_comp\_injective} \\
Theorem~\ref{thm:pxmem} & reduction & \texttt{reach\_le}, \texttt{output\_le}, \texttt{search\_guarantee}: non-worsening replacements keep a member no worse than the seed; the Pivot guarantee~\citep{AilonCN08} is a hypothesis \\
\bottomrule
\end{tabular}
\end{table}

The formal statements connect to the code in two places.  The bounds we
report are values of the expression that \texttt{far\_dual\_bound} bounds,
evaluated by the checker of Section~\ref{sec:cert}.  The move, swap and
contraction formulas are tested against recomputation: exact costs of random
clusterings on contracted instances, and the tracked cost against a full
recomputation at the end of annealing runs with many moves, swaps and
rollbacks.  These tests check final states, not every step.

\section{Development log}\label{app:devlog}

This appendix records which data informed which decision, so that the
reader can judge which results are out of sample.

\begin{itemize}
\item \emph{Development data.}  The fifteen SNAP graphs of
Table~\ref{tab:instances} marked ``dev'' and the 200 PACE~2021 exact-track
instances.  All parameters of the bound (block size $2500$, $|H|\le8$ for local
subgraph rows, ruin regions of about $12$ vertices, half of the budget for the
bound), of the LNS ($|B|\le60$) and of PXMem (population $4$, regions of at
most $30$ nodes, bandit settings, the degree bias of the region centres)
were chosen by looking at runs on these instances.  Several PXMem components
were added after inspecting where KaPoCE's solutions were better than ours on
development graphs (Section~\ref{sec:primal}).
\item \emph{Earlier comparisons.}  A head-to-head comparison with KaPoCE on
the development graphs and on the PACE heuristic track was carried out before
the held-out set was fixed.  Its raw results remain in the repository; it is
not the basis of any claim in this paper, because the same graphs were used
for tuning and because it measured the two solvers asymmetrically (KaPoCE's
time included reading the input, PXMem's did not, and the order was fixed).
\item \emph{Machines.}  Earlier runs of the PACE exact track on a 4-core
workstation with four runs in parallel, also with checked certificates,
proved optimality on 111 of the 173 solved instances.  All numbers in this
paper are from the runs on the Colab machines described in
Section~\ref{sec:exp-setup}, which give \numPaceOptOurs{}; we report the
latter because the other bounds of Table~\ref{tab:pace-bounds} were computed
on the same machines and because their certificates are the archived ones.
\item \emph{Held-out data.}  The twelve SNAP graphs marked ``held-out'' in
Table~\ref{tab:instances} and the analysis plan of
Section~\ref{sec:exp-setup} were fixed and committed to the repository
(\texttt{experiments/heldout.md}) before any run of either solver on them.
No parameter was changed afterwards.  The PACE heuristic-track instances with
the same vertex count as a development graph (ten instances) are excluded
from the held-out PACE summary.
\end{itemize}
